% Image credits for the photographs used in this volume (Book 1).
% Machine-readable license records live in images/book1/CREDITS.md; keep
% the two files in sync when adding or removing a photograph.
\chapter*{Kredit Gambar}

Gambar-gambar dalam buku ini adalah karya asli. Foto-foto dipakai dengan
lisensi bebasnya masing-masing, dengan terima kasih kepada para
pembuatnya:

\begin{itemize}
  \item \emph{John Dalton} (bab atom dan molekul): gravir karya Charles
        Turner berdasarkan Joseph Allen, 1834, Library of Congress
        melalui Wikimedia Commons, domain publik.
  \item \emph{Tabel unsur Dalton, 1808} (bab atom dan molekul): lempeng
        dari buku John Dalton \emph{A New System of Chemical
        Philosophy}, hasil pindaian Internet Archive melalui Wikimedia
        Commons, domain publik.
  \item \emph{Bauksit} (bab bahan baku): foto oleh Werner Schellmann
        melalui Wikimedia Commons, CC\ BY-SA\ 2.5.
  \item \emph{Granit} (bab zat murni): foto oleh Eurico Zimbres melalui
        Wikimedia Commons, CC\ BY-SA\ 2.0\ br.
  \item \emph{Tembaga sulfat anhidrat} (bab mengenali zat): foto oleh
        W.~Oelen melalui Wikimedia Commons, CC\ BY-SA\ 3.0.
  \item \emph{Kristal tembaga sulfat} (bab mengenali zat): foto oleh
        Stephanb melalui Wikimedia Commons, CC\ BY-SA\ 3.0.
  \item \emph{Antoine-Laurent Lavoisier dan Marie-Anne Paulze Lavoisier}
        (bab persamaan setara): lukisan karya Jacques-Louis David, 1788,
        Metropolitan Museum of Art melalui Wikimedia Commons, domain
        publik.
  \item \emph{Telepon bakelit} (bab plastik): foto oleh Holger Ellgaard
        melalui Wikimedia Commons, CC\ BY-SA\ 3.0.
  \item \emph{Ernest Rutherford} (bab isi atom): foto, George Grantham
        Bain Collection, Library of Congress, melalui Wikimedia Commons,
        domain publik.
  \item \emph{Halit} (bab ion): foto oleh Robert M.~Lavinsky melalui
        Wikimedia Commons, CC\ BY-SA\ 3.0.
  \item \emph{Patung Liberty} (bab logam dan korosi): foto oleh
        Elcobbola melalui Wikimedia Commons, domain publik.
  \item \emph{Dmitri Mendeleev} dan \emph{tabelnya tahun 1869} (bab
        tabel periodik): melalui Wikimedia Commons, domain publik.
  \item \emph{Nebula Kepiting} (bab unsur-unsur): NASA, ESA, J.~Hester,
        dan A.~Loll, domain publik.
  \item \emph{Natrium} (bab kulit elektron): foto oleh Dnn87 melalui
        Wikimedia Commons, CC\ BY-SA\ 3.0.
  \item \emph{Klorin di dalam ampul} (bab kulit elektron): foto oleh
        W.~Oelen melalui Wikimedia Commons, CC\ BY-SA\ 3.0.
  \item \emph{Amedeo Avogadro} (bab mol): litografi berdasarkan gambar
        karya C.~Sentier, 1856, koleksi Edgar Fahs Smith, melalui
        Wikimedia Commons, domain publik.
  \item \emph{Kristal salju} (bab kepolaran): foto karya Wilson Bentley,
        1902, melalui Wikimedia Commons, domain publik.
  \item \emph{Louis Pasteur} (bab stereokimia): foto oleh Paul Nadar,
        melalui Wikimedia Commons, domain publik.
  \item \emph{Kristal amonium tartrat} (bab stereokimia): foto oleh
        Marie-Lan Taÿ Pamart melalui Wikimedia Commons, CC\ BY\ 4.0.
  \item \emph{Svante Arrhenius} (bab asam dan basa): fotogravur, 1909,
        melalui Wikimedia Commons, domain publik.
  \item \emph{Tumpukan Volta} (bab sel elektrokimia): foto oleh GuidoB
        melalui Wikimedia Commons, CC\ BY-SA\ 3.0.
  \item \emph{Wallace Carothers di laboratoriumnya} (bab polimer):
        fotografer tidak diketahui, melalui Wikimedia Commons, domain
        publik.
\end{itemize}

Berkas \texttt{images/book1/CREDITS.md} di repositori buku ini
mencantumkan tautan sumber dan URL lisensi lengkap untuk setiap foto.
Piktogram
bahaya adalah piktogram Sistem Harmonisasi Global untuk Klasifikasi dan
Pelabelan Bahan Kimia (GHS), sebagaimana digambar untuk paket
\texttt{ghsystem} dari \TeX\ Live.

\bigskip
Selain foto-foto dan gambar-gambar asli, sebagian gambar adalah
ilustrasi yang dibuat dengan kecerdasan buatan, dihasilkan dengan
pembangkit gambar OpenAI dari perintah yang ditulis oleh para penyunting
buku ini, dan diperiksa ketepatan kimianya sebelum dimuat. Perintah
lengkap untuk setiap gambar yang dihasilkan tercatat dalam
\texttt{images/book1/ai/PROMPTS.md} di repositori.
\clearpage
